\section{Introduction}
Indonesian customer reviews and online comments contain abbreviations, spelling variants, and morphological transformations that differ from standard written forms~\cite{salsabila2018,wibowo2021}. Sentiment classifiers can receive this wording directly or replace known informal forms with dictionary alternatives. Substitution may reduce variation, but ambiguous mappings and context-dependent meanings can also change the input inappropriately.

Prior work examines normalization with Indonesian pretrained models~\cite{bustamin2025,budianoor2026} and compares broader preprocessing pipelines across representations~\cite{prichanti2026}. Our contribution is the combined design: one isolated dictionary intervention, shared audited splits for two public datasets, separate word and character TF-IDF models alongside IndoBERT, and both retrained-pipeline and fixed-model input comparisons. None of these ingredients alone constitutes a new model or a first normalization experiment.

Three questions guide the study. \emph{RQ1:} Does normalization benefit the word model more than the character model or IndoBERT? \emph{RQ2:} Are the effects consistent across PRDECT-ID and SmSA? \emph{RQ3:} Which predictions change when the same trained classifier receives normalized rather than original input? Direct contrasts address RQ1; dataset-specific effects address RQ2; and paired input-switch transitions address RQ3. Absolute scores across datasets with different label schemes are not treated as a controlled comparison of dataset difficulty. Conclusions concern the tested resources and configurations, not cross-domain transfer or semantic preservation.
